\subchapter{Эксперименты по оптимальному размещению объектов обслуживания с учетом изменения транспортной}
\label{ch:optimal-placement-local}

Для того чтобы проверить возможность реализации оптимизации размещения объектов спроса с помощью изменений в транспортной сети, необходимо провести следующие эксперименты:

\begin{itemize}
    \item Провести эксперимент по оценке возможного влияния изменения транспортной доступности к объектам обслуживания за счет существующей сети ОТ;
    \item Провести моделирование потенциального улучшения транспортной доступности между кварталами за счет улучшения сети ОТ на примере небольшого участка города, в задаче оптимального размещения;
    \item Провести сравнение при изменении дорожной сети (добавлении дороги) и изменений транспортной сети (генерации маршрута) на реальном сценарном проекте -- для оценки различий с уже существующим методом;
    \item Провести сравнительный анализ различных целевых параметров связности территории при генерации маршрутов ОТ.
\end{itemize}

\section{Постановка экспериментальной проверки}
\label{sec:placement-experiment-setup}

В данной экспериментальной части рассматривается вопрос непосредственного обеспечения заданного уровня транспортной доступности объектов спроса. В данной работе рассматривается доступность на общественном транспорте или пешком. Использование личного моторизированного или немоторизированного транспорта в данной работе не рассматривается, поскольку учет этих видов передвижения значительно влияет на постановку и рассмотрение задачи, а также вводит дополнительные ограничения. Как было указано ранее в Главе 1, доступность может быть рассчитана с или без учета вместимости объектов обслуживания. В данной работе время в пути понимается как постоянная во времени величина. Изменения времени в пути в зависимости от времени суток / дня недели в этой работе не рассматривается, так как эти уточнения относятся к другому классу оптимизационных задач и не являются необходимыми для экспериментов, представленных в данной работе.

Доступность до объекта обслуживания на общественном транспорте состоит из нескольких этапов: путь от дома до остановки, от остановки у дома до остановки у сервиса (с или без учета пересадок), путь от остановки до места назначения. Чтобы обеспечить суммарное время в пути не более заданного значения (в данном случае -- не более 15 минут), необходимо рассмотреть возможности улучшения каждого участка пути.

Так, время в пути от и до остановок может быть улучшено либо за счет появления новых пешеходных зон/проходов, либо за счет добавления или переноса непосредственно остановки (без учета других подвозящих опций). Оптимизация времени в пути между остановками, на общественном транспорте, находится в фокусе данной части экспериментов, наравне с оптимальным размещением объектов обслуживания (что, непосредственно, уменьшает весь путь между точками спроса и сервисом).





%\section{Базовая оценка транспортно-сервисной доступности}
%\label{sec:baseline-pt-service-access}
%
%\subsection{Декомпозиция времени доступа}
%\label{subsec:access-time-decomposition}

Для каждой пары ``квартал спроса -- квартал сервиса'' полное время доступа задается как
\[
T_{is}=T^{\text{дом}}_{i,p}+T^{\text{ОТ}}_{p,q}+T^{\text{сервис}}_{q,s},
\]
где $T^{\text{дом}}_{i,p}$ --- время пешего подхода от зоны спроса $i$ к остановке $p$, $T^{\text{ОТ}}_{p,q}$ --- время поездки между остановками $p$ и $q$ по сети общественного транспорта, а $T^{\text{сервис}}_{q,s}$ --- время пешего подхода от остановки $q$ к сервисному объекту $s$. Доступ считается обеспеченным, если существует хотя бы один сервис $s$ данного типа, для которого $T_{is}$ не превышает нормативные значения.

Если $T^{\text{дом}}_{i,p}$ превышает допустимое значение, квартал находится далеко от входа в транспортную сеть. Если основной вклад во времени в пути дает $T^{\text{ОТ}}_{p,q}$, ограничение связано с конфигурацией маршрутной сети общественного транспорта или временем пересадок. В случае, если значение $T^{\text{сервис}}_{q,s}$ само по себе больше заданного порога на всё время поездки, квартал, в котором находится сервис, формально достижим транспортом, но плохо связан с ближайшими остановками.

%\subsection{Показатели базового состояния}
%\label{subsec:baseline-access-indicators}
%
%Для базового сценария оценки доступности рассчитываются следующие показатели:
%\begin{itemize}
%    \item доля населения и жилых кварталов, для которых ближайшая остановка находится дальше заданного пешего порога;
%    \item доля населения и жилых кварталов, для которых сервис доступен пешком от остановки назначения только с превышением порога;
%    \item доля населения и жилых блоков, для которых полная цепочка доступа превышает общий временной норматив;
%    \item разложение недоступного спроса по доминирующей причине: участок до остановки, участок поездки на общественном транспорте, участок от остановки до сервиса или нехватка емкости сервиса;
%    \item величина практического неудовлетворенного спроса, рассчитанная после учета емкости доступных сервисов.
%\end{itemize}

Отдельно фиксируется различие между спросом, который не имеет достижимого сервиса вообще, и спросом, который имеет достижимый сервис, но не может быть обслужен из-за ограниченной емкости, поскольку неудовлетворенный спрос может оставаться после распределения спроса по объектам только по значению доступности.

Для проверки того, какая часть цепочки доступа ограничивает достижимость конкретного сервиса, для каждого жилого здания была выполнена классификация результата доступа относительно порога 15 минут. На рисунках \ref{fig:chapter4-polyclinic-access-diagnostics} и \ref{fig:chapter4-school-access-diagnostics} показаны карты этой диагностики для поликлиник и школ по рассматриваемым городам.

\begin{figure}[p]
    \centering
    \includegraphics[width=\linewidth]{ch4/access_diagnostics/polyclinics_access_diagnostics_ru}
    \caption{Диагностика доступности от жилых зданий до поликлиник для 20 городов в доступности 15 минут.}
    \label{fig:chapter4-polyclinic-access-diagnostics}
\end{figure}

\begin{figure}[p]
    \centering
    \includegraphics[width=\linewidth]{ch4/access_diagnostics/schools_access_diagnostics_ru}
    \caption{Диагностика доступности от жилых зданий до школ для 20 городов в доступности 15 минут.}
    \label{fig:chapter4-school-access-diagnostics}
\end{figure}

\section{Эксперимент с изменением матрицы транспортной достижимости}
\label{sec:icccsa-connectivity-placement-experiment}

Исходная проверка совместной логики транспортной связности и размещения сервисов была выполнена в работе \cite{kontsevik2024connectivityPlacement} и опирается на трактовку транспортной связности городских территорий через интермодальную доступность между городскими блоками \cite{morozov2023transportConnectivity}.
\subsection{Гипотеза и постановка задачи покрытия}
\label{subsec:icccsa-hypothesis-location-coverage}

В эксперименте исследовалась следующая гипотеза: улучшение связности между кварталами не более чем на 40\% может уменьшить минимальное число новых поликлиник, оптимально размещенных на территории (в данном случае, это не более 10 минут).

Качество жизни человека в этой постановке рассматривалось как средняя оценка относительно обеспеченности различными сервисами и оценки пространственного несоответствия между жильем и работой. Для экспериментальных целей рассматривался сервис ``поликлиника'', расположение которого оптимизировалось. Как было показано в обзоре существующих решений, алгоритм MCLP дает решения для максимизации покрытия спроса, то есть обеспеченности всех людей, проживающих на заданной территории. Используется улучшенная версия алгоритма MCLP --- CLSCP-SO (Capacitated Location Set Covering Problem-System Optimal), которая учитывает спрос и емкость сервисов. Однако этот алгоритм работает с матрицей стоимости как с константой. В данном случае матрица стоимости является матрицей доступности между выбранными кварталами, которая оптимизируется относительно числа размещаемых поликлиник.

\subsection{Адаптация генетического алгоритма}
\label{subsec:icccsa-genetic-adaptation}

В работе рассматривается изменение значений матрицы стоимости, то есть времени доступности между выбранными блоками, для того чтобы уменьшить требуемое число сервисов для покрытия неудовлетворенного спроса. Гибкость разработанного метода заключается в возможности настраивать предел улучшения доступности сервисов и учитывать любое распределение величин спроса по блокам.

В качестве примера поликлиника рассматривалась как один из жизненно важных сервисов, напрямую влияющих на качество жизни жителей. Для разных сервисов, обеспечивающих базовые потребности, существуют нормативы доступности. В данном случае для сервиса ``поликлиника'' этот норматив составляет 15 минут. Следовательно, если сервис не попадает в зону доступности жилого дома, его жители считаются необслуженными.

Для улучшения обеспеченности жителей необходимым сервисом нужно сделать его доступным: улучшить доступность существующих сервисов и/или добавить новые сервисы. Поскольку в работе рассматривается доступность сервиса на общественном транспорте, которую городские управленцы могут улучшать через реформы, предлагается немного улучшить транспортную доступность между выбранными кварталами.

Так как норматив для выбранного сервиса ``поликлиника'' составляет 15 минут, а выбранный предел улучшения доступности --- до 40\%, необходимо рассматривать те связи между кварталами, которые больше установленного норматива, то есть больше 15 минут, но при улучшении связности могут попасть в заданный норматив. Поэтому необходимо рассматривать связи между кварталами, значение которых не превышает 25 минут, поскольку $25 \cdot 0.6 = 15$.
%
%Пример матрицы стоимости приведен в таблице \ref{tab:chapter4-tikhevich-cost-matrix}. Цветом отмечены значения, рассматриваемые для оптимизации.

%\begin{table}[p]
%    \centering
%    \scriptsize
%    \caption{Пример матрицы стоимости из эксперимента}
%    \label{tab:chapter4-tikhevich-cost-matrix}
%    \resizebox{\linewidth}{!}{%
%    \begin{tabular}{|c|c|c|c|c|c|c|c|c|c|c|}
%        \hline
%        \textbf{идентификатор блока} & \textbf{0} & \textbf{1} & \textbf{2} & \textbf{3} & \textbf{4} & \textbf{5} & \textbf{6} & \textbf{7} & \textbf{8} & \textbf{9} \\
%        \hline
%        0 & 0.0 & 11.3 & \cellcolor{blue!25}21.0 & \cellcolor{blue!25}25.0 & 27.6 & 26.5 & 34.6 & 32.4 & 26.1 & \cellcolor{blue!25}25.0 \\
%        1 & 11.3 & 0.0 & \cellcolor{blue!25}24.5 & 27.3 & 29.9 & 28.8 & 36.9 & 34.7 & 28.4 & 27.3 \\
%        2 & \cellcolor{blue!25}17.3 & \cellcolor{blue!25}20.4 & 0.0 & 13.8 & \cellcolor{blue!25}22.8 & \cellcolor{blue!25}21.8 & 29.9 & 27.7 & \cellcolor{blue!25}21.3 & \cellcolor{blue!25}20.2 \\
%        3 & 25.9 & 29.0 & 13.8 & 0.0 & \cellcolor{blue!25}17.2 & 29.0 & 37.3 & 35.1 & \cellcolor{blue!25}20.3 & \cellcolor{blue!25}19.2 \\
%        4 & 26.3 & 29.4 & 26.6 & \cellcolor{blue!25}18.6 & 0.0 & \cellcolor{blue!25}21.2 & 35.6 & 33.4 & 25.8 & \cellcolor{blue!25}24.4 \\
%        \hline
%    \end{tabular}%
%    }
%\end{table}

В данном случае генетический алгоритм используется в стандартной реализации. Он работает, имитируя процесс естественного отбора для оптимизации решений сложных задач. В основе генетического алгоритма лежит начальный набор потенциальных решений, представленных как исходная популяция. Затем эти решения оцениваются на основе функции приспособленности, которая определяет, насколько хорошо они выполняют задачу. Через процессы отбора, скрещивания и мутации алгоритм итеративно развивает популяцию, способствуя воспроизводству индивидов с большей приспособленностью. На протяжении нескольких поколений популяция стремится к оптимальным или почти оптимальным решениям, что делает генетические алгоритмы особенно эффективными для решения оптимизационных задач в различных областях, включая городские исследования.

Важно отметить, как именно задаются начальная популяция и функция приспособленности. В качестве начальной популяции выбираются матрицы с измененными рассматриваемыми значениями:
\[
M_i \in \mathbb{R}^{n \times n} \quad \text{для } i = 1, 2, \ldots, \text{population\_size},
\]
\[
M_{i,ij} =
\text{random\_uniform}(0.6, 1) \times \text{cost\_matrix}_{ij},
\quad \text{если } (i,j)\in\text{selected cells}.
\]

В этом случае результат алгоритма оптимального размещения сервисов, то есть число размещенных сервисов, используется как функция приспособленности:
\[
F(M) = \text{count}(\text{clscpso.fac2cli}[M]).
\]

В результате алгоритм выбирает лучшее решение, содержащее минимальное число новых сервисов, которые необходимо разместить.

\subsection{Территория исследования и данные}
\label{subsec:icccsa-study-area-data}

В работе рассматривается сценарий размещения нового промышленного предприятия в городе в рамках развития приоритетных промышленных секторов согласно стратегии социально-экономического развития города. Рассматривается размещение промышленного предприятия в портово-промышленном районе Санкт-Петербурга. Размещаемое предприятие и выбранные городские блоки представлены на рисунке \ref{fig:chapter4-tikhevich-site-selection}.

\begin{figure}[p]
    \centering
    \includegraphics[width=\linewidth]{ch4/optimal_local/tikhevich_site_selection}
    \caption{Выбранная территория эксперимента.}
    \label{fig:chapter4-tikhevich-site-selection}
\end{figure}

Рассматриваемые кварталы выбирались следующим образом: согласно гипотезе пространственного несоответствия между местом работы и жильем, необходим баланс между заработной платой работника, потенциальной средней арендной платой и расстоянием от работы. Заработная плата потенциальных работников была взята как средняя для квалифицированных промышленных работников города.

Все необходимые данные о размещении сервисов, сетях общественного транспорта и параметрах жилых зданий, необходимых для расчета оценочной численности населения, были собраны из OpenStreetMap. Данные о мощности сервисов и спросе рассчитывались согласно работе \cite{khrulkov2022city}. Реализованный алгоритм опирается на гравитационную модель, ограниченное число жителей и ограниченные мощности сервисов.

\subsection{Реализация метода и результаты}
\label{subsec:icccsa-results}

Текущая оценка доступности ``на рассматриваемый момент времени'' была выполнена с помощью библиотеки BlocksNet \cite{blocks}. Интермодальный граф для выбранной территории и оценка доступности жителей выбранных кварталов к сервису ``поликлиника'' представлены на рисунке \ref{fig:chapter4-tikhevich-input-data}. Каждый сервис имеет собственный норматив доступности согласно ГОСТ. Для поликлиники норматив доступности составляет 15 минут.

\begin{figure}[p]
    \centering
    \begin{minipage}[t]{0.49\linewidth}
        \centering
        \includegraphics[width=\linewidth]{ch4/optimal_local/tikhevich_intermodal_graph}
        \smallskip
        \small а) Геослой с интермодальным графом общественного транспорта
    \end{minipage}\hfill
    \begin{minipage}[t]{0.49\linewidth}
        \centering
        \includegraphics[width=\linewidth]{ch4/optimal_local/tikhevich_initial_state}
        \smallskip
        \small б) Блоки, покрытые и не покрытые сервисом
    \end{minipage}
    \caption{Входные данные эксперимента.}
    \label{fig:chapter4-tikhevich-input-data}
\end{figure}

В работе все новые поликлиники имеют фиксированную среднюю мощность 400 человек. Следует отметить, что необходимым, но недостаточным условием для того, чтобы считать квартал обеспеченным сервисом, является его попадание в зону доступности сервиса. Доступность рассчитывалась между кварталами на интермодальном графе, построенном на открытых данных OpenStreetMap.

\begin{figure}[p]
    \centering
    \begin{minipage}[t]{0.49\linewidth}
        \centering
        \includegraphics[width=\linewidth]{ch4/optimal_local/tikhevich_ga_fitness}
        \smallskip
        \small а) Прогресс функции приспособленности в генетическом алгоритме
    \end{minipage}\hfill
    \begin{minipage}[t]{0.49\linewidth}
        \centering
        \includegraphics[width=\linewidth]{ch4/optimal_local/tikhevich_accessibility_delta_distribution}
        \smallskip
        \small б) Распределение требуемых улучшений доступности, мин.
    \end{minipage}
    \caption{Результаты генетического алгоритма.}
    \label{fig:chapter4-tikhevich-ga-results}
\end{figure}

Так как число возможных комбинаций в данном случае велико, генетический алгоритм был выбран как инструмент для поиска субоптимального решения. На вход генетический алгоритм получает матрицу связности и изменяет значения только в заданных ячейках, то есть кварталы, которые не удовлетворяют условию для оптимизации, не изменяют связность между собой. В результате алгоритм CLSCP-SO рассчитывает оптимальное число новых поликлиник, необходимых для удовлетворения спроса. Результаты алгоритма представлены на рисунке \ref{fig:chapter4-tikhevich-ga-results}.

Рисунок \ref{fig:chapter4-tikhevich-ga-results} показывает, что за счет оптимизации времени доступности между блоками требуемое число новых сервисов было уменьшено с 13 до 9. Распределение изменения транспортной связности между кварталами показывает, что среднее время, необходимое для улучшения транспортной связности, составляет 10 минут.

\begin{figure}[p]
    \centering
    \begin{minipage}[t]{0.49\linewidth}
        \centering
        \includegraphics[width=\linewidth]{ch4/optimal_local/tikhevich_new_services}
        \smallskip
        \small а) Размещение существующих и дополнительных сервисов
    \end{minipage}\hfill
    \begin{minipage}[t]{0.49\linewidth}
        \centering
        \includegraphics[width=\linewidth]{ch4/optimal_local/tikhevich_link_improvements}
        \smallskip
        \small б) Распределение изменения транспортной связности между кварталами
    \end{minipage}
    \caption{Итоговое размещение сервисов и улучшения транспортной связности на рассматриваемой территории.}
    \label{fig:chapter4-tikhevich-services-and-links}
\end{figure}

Рисунок \ref{fig:chapter4-tikhevich-services-and-links} показывает блоки, в которых должен быть размещен сервис: сервис обозначен зеленым маркером; блоки, в которых сервис уже существовал ранее, обозначены красным маркером; блоки, не покрытые до оптимизации и размещения дополнительных сервисов, обозначены красным цветом. На этом же рисунке показано уменьшение времени перемещения между блоками: линии окрашены градиентом в зависимости от требуемой величины изменения.

Можно отметить, что в некоторых случаях новый сервис размещается в кварталах, где сервис уже существует. Это указывает на необходимость расширения или оптимизации существующего сервиса. Также можно наблюдать, что из-за характера маршрутов и дорожной сети в некоторых случаях требуется только небольшое улучшение связности между удаленными кварталами, от 2 до 4 минут.

Таким образом, предложенные алгоритмом улучшения связей между кварталами следует интерпретировать как предложения по улучшению транспортной связности территории. При этом работа не задает конкретное проектное решение в виде нового маршрута, расписания или дорожного объекта. Ее результат состоит в том, что при улучшении связей между территориями, даже без явного добавления новых дорог, можно уменьшить число необходимых сервисных объектов.

\section{Эксперимент с добавлением маршрута общественного транспорта}
\label{sec:add-route-experiment}

Сценарий добавления маршрута используется для проверки того, может ли транспортная интервенция снизить недоступность без открытия новых сервисных объектов. В эксперименте набор сервисов и их емкости сохраняются неизменными, а в транспортную сеть добавляется новая линия или связь между существующими остановками. После этого пересчитываются кратчайшие времена $T^{\text{ОТ}}_{p,q}$ и полная доступность $T_{is}$.

Основные выходные показатели сценария:
\begin{itemize}
    \item изменение доли населения с доступом к сервису в пределах норматива;
    \item изменение суммарного практического неудовлетворенного спроса;
    \item число блоков, для которых доминирующая причина провала изменилась после добавления маршрута;
    \item распределение выигрыша по времени доступа;
    \item список зон, для которых маршрут создает новый достижимый сервисный выбор.
\end{itemize}





\section{Сравнение добавления дороги и маршрута на примере Санкт-Петербурга}
\label{sec:spb-road-vs-route}

Для Санкт-Петербурга рассматриваются три сценария поверх одного базового состояния:
\begin{enumerate}
    \item добавление дорожной связи;
    \item добавление маршрута общественного транспорта;
    \item совместное добавление дорожной связи и маршрута.
\end{enumerate}




\begin{figure}[p]
    \centering
    \includegraphics[width=\linewidth]{ch4/optimal_local/telmana_connector_context_row}
    \caption{Исходные слои эксперимента для квартала Тельмана: проект планировки, дорожная сеть и кварталы, существующие маршруты общественного транспорта, существующие поликлиники и базовый неудовлетворенный спрос.}
    \label{fig:spb-telmana-connector-context}
\end{figure}

\begin{figure}[p]
    \centering
    \includegraphics[width=\linewidth]{ch4/optimal_local/telmana_connector_scenarios_worse_new_services}
    \caption{Сценарии для квартала Тельмана с большим числом новых сервисов. Во всех четырех вариантах требуется открыть 14 новых поликлиник; сценарии приведены перед группой лучших вариантов, чтобы сначала описать случаи с большей потребностью в сервисной интервенции.}
    \label{fig:spb-telmana-connector-worse-new-services}
\end{figure}

\begin{figure}[p]
    \centering
    \includegraphics[width=\linewidth]{ch4/optimal_local/telmana_connector_scenarios_better_new_services}
    \caption{Сценарии для квартала Тельмана с меньшим числом новых сервисов. Добавление маршрута общественного транспорта снижает требуемое число новых поликлиник c 14 до 13 в сценариях с проектным кварталом и до 12 в сценариях без него; длина сгенерированного маршрута составляет 10 остановок.}
    \label{fig:spb-telmana-connector-better-new-services}
\end{figure}

\begin{figure}[p]
    \centering
    \includegraphics[width=\linewidth,height=0.82\textheight,keepaspectratio]{ch4/optimal_local/polyclinic_route_strategy_7x3_ru}
    \caption{Сравнение трех стратегий генерации маршрутов общественного транспорта для снижения потребности в новых поликлиниках. Зеленая рамка отмечает сценарий с минимальным числом новых объектов обслуживания.}
    \label{fig:chapter4-polyclinic-route-strategy-scenarios}
\end{figure}
